\documentclass[aspectratio=169]{beamer}
\usetheme{Madrid}
\usecolortheme{default}
\usepackage[T1]{fontenc}
\usepackage[utf8]{inputenc}
\usepackage[spanish,es-nodecimaldot]{babel}
\usepackage{amsmath}
\usepackage{booktabs}
\usepackage{graphicx}
\usepackage{siunitx}

\title{Laboratorio 1}
\subtitle{Python científico reproducible}
\author{Cristian Stiven Capera, Ana Sofia García, Jeimmy Andrea Gonzáles, Karol Mariana Gutiérrez y Brayan Andrés Ruiz}
\date{\today}

\begin{document}

% Diapositiva 1: Portada
\begin{frame}
    \titlepage
\end{frame}

% Diapositiva 2: Problema y arquitectura (Secciones 7.1 y 7.4)
\begin{frame}{Problema y arquitectura}
    \begin{block}{Pregunta de Trabajo y Enfoque}
        ¿Puede el equipo reconstruir en un computador distinto un análisis reproducible y legible, entendiendo que como ingenieros no siempre tendremos el privilegio de usar nuestros propios equipos?
    \end{block}
    \vspace{-0.1cm}
    \begin{itemize}
        \small
        \item \textbf{Reproducibilidad:} Garantizar que el código corra en cualquier equipo de computación externo sin requerir copias manuales del entorno virtual ni ajustes de rutas.
        \item \textbf{Legibilidad y Estándares:} Escribir código limpio y estructurado que cualquier persona pueda leer, evaluar y mantener fácilmente.
        \item \textbf{Arquitectura Modular:} Separación de responsabilidades en carpetas (\texttt{src/}, \texttt{tests/}, \texttt{scripts/}) controladas mediante \texttt{uv}.
    \end{itemize}
\end{frame}

% Diapositiva 3: Métodos y validación (Secciones 7.5, 7.6, 7.7 y 7.10)
\begin{frame}{Métodos y validación}
    \begin{columns}[T]
        \column{0.5\textwidth}
        \textbf{Matemática y Simulación:}
        \begin{itemize}
            \small
            \item Semilla determinista $\text{seed} = 42$ y $\Delta t = 0.05\text{ s}$.
            \item Modelo de señal cruda: $x_{\text{raw}}(t) = \sin(2\pi \cdot 0.5 t) + 0.2 \cdot \mathcal{N}(0,1)$
            \begin{itemize}
                \tiny
                \item Onda senoidal de frecuencia de 0.5Hz y ruido gaussiano.
            \end{itemize}
            \item Variante centrada de la señal cruda, se le resta la media temporal.
            \item Estadística descriptiva, RMS y desviación estándar con $ddof=1$.
            \item Matriz de Mezcla ($\mathbf{A}$) bien condicionada:
            $\mathbf{A} = \begin{pmatrix} 2.0 & 1.0 \\ 1.0 & 3.0 \end{pmatrix}$.
        \end{itemize}
        
        \column{0.5\textwidth}
        \textbf{Control de Calidad y Pruebas:}
        \begin{itemize}
            \small
            \item Validación cruzando listas y conjuntos: analizando \texttt{observed = [SYNTH\_A, SYNTH\_B, SYNTH\_A]} frente a requeridos \texttt{\{SYNTH\_A, SYNTH\_B, REFERENCE\}}.
            \begin{itemize}
                \tiny
                \item Duplicados: Detectó \texttt{SYNTH\_A}.
                \item Faltantes: Ausencia de \texttt{REFERENCE}.
                \item Rechazo de nombres vacíos.
            \end{itemize}
            \item Derivada e integración numérica:
            \begin{itemize}
                \tiny
                \item Dev: Diferencias finitas para evaluar sensibilidad al ruido.
                \item Int: Validación trapezoidal con función $t^2$.
            \end{itemize}
            \item Pruebas unitarias exitosas con \texttt{pytest}.
        \end{itemize}
    \end{columns}
\end{frame}

% Diapositiva 4: Resultados y Evidencia.
\begin{frame}{Resultados y Evidencia Generada}
    \begin{columns}[T]
        \column{0.45\textwidth}
        \centering
        \textbf{Visualización del Pipeline}
        \vspace{0.05cm}
        \includegraphics[width=0.70\linewidth]{figures/summary.png}
        \tiny
        \begin{itemize}
            \item \textbf{Panel A/B/C:} Comparativa de señal cruda y centrada, Histograma de amplitudes de la señal centrada, y Derivada numérica en el intervalo $[1.0, 2.0]$ segs.
        \end{itemize}
        
        \column{0.55\textwidth}
        \textbf{Artefactos Dinámicos (Generados por \texttt{reproduce.py}):}
        \vspace{0.05cm}
        \begin{itemize}
            \small
            \item \textbf{\texttt{results/stats.json}:}
            \begin{itemize}
                \tiny
                \item Estadísticas descriptivas de la señal cruda y centrada con unidades en mV ($ddof=1$).
            \end{itemize}
            \item \textbf{\texttt{results/summary\_table.tex}:}
            \begin{itemize}
                \tiny
                \item Tabla LaTeX insertada de forma dinámica en el informe sin transcripción manual.
            \end{itemize}
            \item \textbf{\texttt{results/quality\_report.json}:}
            \begin{itemize}
                \tiny
                \item Reporte de validación de canales y rechazo de nombres vacíos.
            \end{itemize}
            \item \textbf{\texttt{results/algebra\_report.json}:}
            \begin{itemize}
                \tiny
                \item Verificación de formas, número de condición y recuperación de canales por álgebra lineal.
            \end{itemize}
            \item \textbf{\texttt{results/derivation\_integration\_report.json}:}
            \begin{itemize}
                \tiny 
                \item Métricas de error (RMS/Máximo) de la derivada y validación de integración.
            \end{itemize}
        \end{itemize}
    \end{columns}
\end{frame}

% Diapositiva 5: Reproducibilidad y limitaciones (Secciones 7.3 y 7.9)
\begin{frame}{Reproducibilidad y limitaciones}
    \begin{itemize}
        \item \textbf{Curva de Aprendizaje y Herramientas:} Retos operativos iniciales en el equipo para dominar la manipulación estricta de entornos con \texttt{uv}, analizar el código fuente de forma automática para detectar errores, problemas de lógica, e inconsistencias con \texttt{Ruff} y automatización vía scripts.
        \item \textbf{Reproducibilidad Lograda:} A pesar de las dificultades técnicas, se consolidó un pipeline capaz de reconstruir el análisis en cualquier equipo mediante comandos limpios.
        \item \textbf{Limitación Analítica y Alcance:} El sistema opera puramente sobre señales sintéticas controladas; no es extrapolable a datos fisiológicos reales.
    \end{itemize}
\end{frame}

\end{document}